\documentclass{article}
\usepackage[utf8]{inputenc}
\usepackage[T1]{fontenc}
\usepackage[english]{babel}
\usepackage[left=3cm, right=3cm, top=2cm, bottom=2cm]{geometry}

\usepackage{newpxtext}
\usepackage{newpxmath}
\usepackage{pgfplots}
\pgfplotsset{compat=1.18}

\usepackage{amsmath}
\usepackage{nccmath}
\usepackage{stmaryrd}
\usepackage{textcomp}
\usepackage{siunitx}
\usepackage{wasysym}
\usepackage{bbm}


\usepackage{tikz}
\usetikzlibrary{decorations.pathreplacing}
\usetikzlibrary{calc,arrows.meta,babel}
\definecolor{planefill}{RGB}{238,240,245}
\definecolor{planeline}{RGB}{201,206,220}
\definecolor{planeedge}{RGB}{160,167,196}
\definecolor{oliveD}{RGB}{138,125,42}
\definecolor{redD}{RGB}{158,43,37}
\definecolor{blueD}{RGB}{31,79,168}
\definecolor{slate}{RGB}{106,111,156}


\usepackage{tabularx}
\usepackage{graphicx}
\usepackage{xcolor}
\usepackage{float}
\usepackage[most,many]{tcolorbox}
\usepackage[hidelinks]{hyperref}

\newcounter{notion}[section]
\renewcommand{\thenotion}{\thesection.\arabic{notion}}
\let\frac\dfrac
\newcommand{\dsum}{\displaystyle\sum}

\definecolor{myblue}{RGB}{60,90,200}
\definecolor{myred}{RGB}{200,60,60}
\definecolor{mygreen}{RGB}{40,120,70}
\definecolor{myorange}{RGB}{190,120,20}

\tcbset{
ideaBox/.style={
    enhanced,
    colback=red!6,
    colframe=myred,
    coltitle=black,
    colbacktitle=red!6,
    fonttitle=\bfseries\scshape,
    attach title to upper={\par\vspace{4pt}\noindent},
    separator sign none,
    boxrule=1.2pt,
    arc=5pt,
    left=10pt, right=10pt, top=8pt, bottom=8pt,
    }
}

\newtcbtheorem[use counter=notion]{definitionbox}{Definition}{ideaBox}{def}
\newtcbtheorem[use counter=notion]{propositionbox}{Proposition}{ideaBox}{prop}
\newtcbtheorem[use counter=notion]{theorembox}{Theorem}{ideaBox}{theo}

% --- Lab-notebook apparatus -------------------------------------------------
% Every experiment in this document is one box, always in the same four beats:
% hypothesis, protocol, measurement, verdict. The frame colour is the verdict,
% so the outcome of a run is readable without reading the run.
\tcbset{
logBox/.style={
    enhanced, breakable,
    coltitle=black,
    fonttitle=\bfseries,
    attach title to upper={\par\vspace{4pt}\noindent},
    separator sign none,
    boxrule=1.0pt,
    arc=4pt,
    left=9pt, right=9pt, top=7pt, bottom=7pt,
    },
failBox/.style={logBox, colback=red!5,    colframe=myred,    colbacktitle=red!5},
winBox/.style={ logBox, colback=green!5,  colframe=mygreen,  colbacktitle=green!5},
openBox/.style={logBox, colback=orange!6, colframe=myorange, colbacktitle=orange!6},
}

\newcounter{exper}
\renewcommand{\theexper}{\arabic{exper}}
\newtcbtheorem[use counter=exper]{expfail}{Experiment}{failBox}{exp}
\newtcbtheorem[use counter=exper]{expwin}{Experiment}{winBox}{exp}
\newtcbtheorem[use counter=exper]{expopen}{Experiment}{openBox}{exp}

\newcommand{\hyp}{\textbf{Hypothesis.\ }}
\newcommand{\proto}{\textbf{Protocol.\ }}
\newcommand{\meas}{\textbf{Measured.\ }}
\newcommand{\verd}{\textbf{Verdict.\ }}

\newcommand{\rejected}{\textcolor{myred}{\textsc{rejected}}}
\newcommand{\adopted}{\textcolor{mygreen}{\textsc{adopted}}}
\newcommand{\untested}{\textcolor{slate}{\textsc{untested}}}
\newcommand{\ongoing}{\textcolor{myorange}{\textsc{running}}}

% Citations as a small raised number, without brackets: [3] inline is heavy in a
% document this dense.
\makeatletter
\renewcommand\@cite[2]{\textsuperscript{#1\ifx\@empty#2\@empty\else\,#2\fi}}
\renewcommand\@biblabel[1]{#1.}
\makeatother

\usepackage{enumitem}
\setcounter{tocdepth}{2}

\usepackage{fancyhdr}
\pagestyle{fancy}
\renewcommand{\headruleskip}{1.5mm}
\renewcommand{\headrulewidth}{0.1pt}
\renewcommand{\footruleskip}{1.5mm}
\renewcommand{\footrulewidth}{0.1pt}

\fancyhead[L]{}
\fancyhead[R]{}
\fancyfoot[L]{}
\fancyfoot[C]{\thepage}
\fancyfoot[R]{}

\title{Raidium Data Challenge 2026}
\author{Alexandre Tranié}
\date{August 2026}

\begin{document}

\maketitle
\tableofcontents
\newpage

\section*{Summary}
\addcontentsline{toc}{section}{Summary}

Segment 54 organs on 2D CT slices, from 800 images whose annotation is
\emph{partial}: on each slice only 27 of the 54 organs were outlined, so a pixel
labelled \texttt{0} means ``background \emph{or} one of the 27 nobody looked
at''. The trajectory:

\begin{center}
\begin{tabularx}{\textwidth}{@{}lXrr@{}}
\hline
\textbf{step} & \textbf{what changed} & \textbf{Dice} & \textbf{silent} \\
\hline
baseline        & U-Net ResNet-34, marginal loss            & 0.4474 & 22 \\
architecture    & DINOv3~\cite{dinov3} + DPT~\cite{dpt} encoder      & 0.4855 & 8 \\
self-training   & pseudo-labels on the 1200 raw slices      & 0.5235 & 6 \\
external data   & TotalSegmentator v1, resliced to 2D       & 0.6803 & 0 \\
+ fine-tuning   & on the 640 challenge slices               & 0.7522 & 0 \\
+ pixel spacing & measured at $1.45$\,mm/px, not $1.20$     & 0.7584 & 2 \\
+ ensemble      & 7 models, 2 architectures, 2 scales       & \textbf{0.7796} & 0 \\
\hline
leaderboard     & (validation $\to$ held-out test)          & \textbf{0.7612} & 0 \\
\hline
\end{tabularx}
\end{center}

\noindent The decisive step was not architectural. It was proving that the 54
anonymous challenge classes \emph{are} TotalSegmentator~v1~\cite{totalseg} with
four merges (\S\ref{sec:identify}), which turns a public dataset of 1204 CT volumes
into complete supervision for exactly this label space. Everything before that
step is worth $+0.076$ in total; that step alone is worth $+0.23$.

\vspace{2.0mm}

\noindent Final standing: \textbf{1\textsuperscript{st}}, $0.0163$ ahead of the
previous leader ($0.7449$), from roughly 21\textsuperscript{st} at the start.


\newpage

\section{Context}

The dataset is composed of 2D CT scan images and their corresponding masks for
different organs (54 in total). The challenge is to develop a model that can
accurately segment the organs from the CT scan images.

\vspace{2.0mm}

Composition of the dataset:
\begin{itemize}
    \item 800 annotated CT scan images in PNG format, $256\times256$ pixels;
    \item 800 corresponding masks;
    \item 1200 CT scan images with no annotation at all;
    \item 500 CT scan images for the test set.
\end{itemize}

\vspace{2.0mm}

\noindent The trick here is that the 800 annotated CT scans are \emph{not fully
annotated}: not every organ is outlined. A label file states, for each image,
which organs were examined - but it only lists 27 of the 54. We therefore lack
any information about the other 27. A pixel set to \texttt{0} does not mean
``background'', it means ``background \emph{or} one of the organs nobody
outlined''.

\vspace{2.0mm}

\noindent Model performance is assessed with the Dice coefficient, a measure of overlap
between two samples. It ranges from 0 to 1, where 1 indicates perfect overlap
and 0 no overlap:
\begin{equation}
    \text{Dice}(A, B) = \frac{2 |A \cap B|}{|A| + |B|}
\end{equation}

\noindent Crucially, the score averages Dice \emph{per organ first}, then over the 54
organs. Every class therefore weighs $1/54$ regardless of how often it appears,
and a class that is never predicted costs $0.0185$ in full.

\newpage

\section{Baseline}
\label{sec:baseline}

A U-Net with a ResNet-34 encoder, trained with the marginal loss~\cite{marginal}
that handles the partial annotation, scores \textbf{0.4474} on our validation split. Two
reference points frame it. Plain cross-entropy - which reads every \texttt{0}
pixel as background and therefore teaches the network to erase the 27 organs it
was not shown - collapses to \textbf{0.1506}. The challenge's own
nnU-Net~\cite{nnunet} baseline sits at \textbf{0.4399}.

\vspace{2.0mm}

\noindent That first gap, $+0.297$ for the loss alone, is the largest single
effect measured anywhere in this project. It is worth stating up front because it
fixes the scale against which everything below should be read: no architectural
change in this document is worth a tenth of it, and the one intervention that
eventually beat it was also about supervision, not about the network.

\subsection{The problem this baseline exposes}

It does not predict the smaller classes at all: 22 of the 54 organs are never
output. Since the metric weighs every class equally, those 22 silent classes cap
the achievable score at $1 - 22/54 = 0.59$ before anything else goes wrong. The
question that drives the rest of the document is therefore not ``how do I segment
better'' but ``how do I stop the model from forgetting two fifths of the label
space''.

\subsection{Warning: there are two ResNet-34 numbers}

$0.4474$ is the best ResNet-34 run of the project, trained to convergence. Several
controlled experiments below use a deliberately shorter recipe (40 epochs, batch
4) whose control arm scores $0.3828$ with 25 silent classes. Both are ``the
ResNet-34 baseline''. They are comparable \emph{within} their own table and never
across tables. Whenever two numbers here seem to contradict each other, check the
training budget first.

\section{The hypothesis log}
\label{sec:log}

Understanding the data is what helped most, and on this challenge that had to come
before any modelling: the identifiers are anatomically anonymous, and ``organ 36
is missing'' is not actionable in the way ``the liver is missing'' is. Naming them
(\S\ref{sec:identify}) is what eventually opened the external data, which is why
it comes first below even though it produces no score by itself.

\vspace{2.0mm}

\noindent The table is the complete list of what was considered, each line with
its verdict. It is meant to be read on its own, as the answer to ``did I already
try this?''. The dead ends take up as much room as the rest of this document,
because ruling out a hypothesis properly costs compute, and writing down
\emph{why} it failed is the only thing that stops it from being tried a second
time.

\begin{center}
\small
\begin{tabularx}{\textwidth}{@{}p{4.3cm}Xl@{}}
\hline
\textbf{hypothesis} & \textbf{what came of it} & \textbf{status} \\
\hline
\multicolumn{3}{@{}l@{}}{\textit{Loss and sampling}}\\
Marginal loss for partial annotation & $0.1506 \to 0.4474$ & \adopted \\
Weight the loss by inverse organ area & $0.4248$, i.e.\ \emph{below} the baseline & \rejected \\
Oversample slices carrying a rare organ & small but real; kept in the final recipe & \adopted \\
Per-class decision thresholds & 54 parameters fitted on 160 images; only a single global threshold at $0.25$ survives cross-validation & \rejected \\
Focal loss, Tversky with $\beta > \alpha$ & only ever run inside the area-weighting arm, never isolated & \untested \\
\hline
\multicolumn{3}{@{}l@{}}{\textit{Architecture}}\\
DINOv3 + DPT instead of ResNet-34 & $+0.038$, the largest architectural gain measured here & \adopted \\
DINOv3 pretrained on the external set & $+0.044$ at pretraining, so the advantage survives. But it takes only $+0.029$ from fine-tuning where ResNet-34 takes $+0.079$, and ends \emph{below} it & \adopted \\
Raise \texttt{encoder-lr-scale} at fine-tuning & monotone and real ($0.1 \to 0.7518$, $1.0 \to 0.7568$) but worth $+0.005$, and \emph{nothing} in the ensemble & \rejected \\
Freeze the DINOv3 encoder & $0.2672$: grey-scale CT is too far from the pretraining distribution & \rejected \\
U-Net with attention, multi-task heads & never reached - the data axis dominated everything & \untested \\
\hline
\multicolumn{3}{@{}l@{}}{\textit{Resolution and field of view}}\\
Upsample the slice to 384 or 512 & non-monotone, and \emph{nothing} on the six smallest organs. Re-run later on a model where everything else works: still $-0.08$, and the adrenal probability moves from $0.006$ to $0.001$ & \rejected \\
Fix the pixel spacing to $1.45$\,mm/px & $+0.006$ after fine-tuning, nothing before it - and it \emph{costs} the two smallest organs, see below & \adopted \\
Zoom cascade: detect, crop, re-segment & $0 \to 0.22$ on the iliac vessels, yet the ablation shows the cause is the narrowed field of view, not the pixel budget & \rejected \\
\hline
\multicolumn{3}{@{}l@{}}{\textit{Input representation}}\\
Recover Hounsfield units, three windows & $0.3059$ against $0.4474$ & \rejected \\
Per-slice contrast handling & useless on the challenge data, then essential when generating the external set (\S\ref{sec:external}) & \adopted \\
\hline
\multicolumn{3}{@{}l@{}}{\textit{Data}}\\
Presence classifier: which organs does this slice contain? & oracle ceiling $+0.007$ --- the model almost never invents an absent organ, so there is nothing to win & \rejected \\
Left--right flip augmentation & $0.4124$ against $0.4474$: it halves the supervision of the classes that have no mirror twin & \rejected \\
Pseudo-label the 1200 unannotated slices & $+0.038$ & \adopted \\
Run TotalSegmentator directly as a labeller & it is a 3D model; on 2D slices stacked into fake volumes its recall falls to 5--30\,\% & \rejected \\
\textbf{Slice annotated 3D volumes into 2D and remap the labels} & \textbf{$+0.23$, and every silent class disappears} (\S\ref{sec:external}) & \adopted \\
\hline
\multicolumn{3}{@{}l@{}}{\textit{Ensembling}}\\
Search the ensemble weights on validation & the winning weighting beats the uniform one by $+0.02$ on those 160 images and by nothing anywhere else & \rejected \\
Average models that differ only in a hyper-parameter & two ResNet-34 fine-tuned at different learning rates: $+0.003$ & \rejected \\
\textbf{Average models that differ in architecture or in pixel scale} & \textbf{$+0.012$ on validation and $+0.016$ on the leaderboard} - the one case where the leaderboard gained \emph{more} than validation predicted & \adopted \\
\hline
\end{tabularx}
\end{center}

\noindent It is the last line that unlocked everything. What follows retraces the
path in the order it was actually walked.

\newpage

\section{Identifying the 54 organs}
\label{sec:identify}

The challenge identifiers are anatomically mute. A first pass matched them
against the 117 classes of TotalSegmentator, by crossing Dice and precision over
127 slices, and named 46 of the 54. Eight resisted: \texttt{7, 8, 12, 39, 41,
48, 52, 53}.

\subsection{The key: the numbering follows alphabetical order}

The decisive fact does not come from the images. Sorting the already-named
identifiers by number, the names come out in alphabetical order - without a
single exception across all 46. That is simply how a \texttt{sorted()} call over
a class dictionary numbers its outputs. Every anonymous identifier is therefore
\emph{bracketed} by its two named neighbours, which shrinks the candidate space
to almost nothing.

\vspace{2.0mm}

\noindent Geometry and boundary adjacency settle the rest:

\begin{center}
\begin{tabularx}{\textwidth}{llX}
\hline
\textbf{id} & \textbf{name} & \textbf{what establishes it} \\
\hline
7, 8 & \texttt{clavicula\_left/right} & elongated bone at the lung apex; 100\,\% of 8's boundary touches \texttt{scapula\_right} \\
12 & \texttt{face} & convex 8{,}700\,px block in front of the face, on skull slices, overflowing the skin \\
39 & \texttt{pancreas} & retro-gastric comma; 46\,\% boundary contact with the splenic vein, 14\,\% duodenum \\
41 & \texttt{pulmonary\_artery} & branching mediastinal trunk at the hilum; 39\,\% contact with the heart, 31\,\% aorta \\
48 & \texttt{spleen} & left upper quadrant; 32\,\% contact with the stomach, 17\,\% left kidney \\
52, 53 & \texttt{vertebrae\_cervical/lumbar} & mutually exclusive with 54, disjoint regional profiles \\
\hline
\end{tabularx}
\end{center}

\subsection{The spine is split into three classes}

\texttt{52}, \texttt{53} and \texttt{54} never coexist: across 190 co-annotated
slices, 52 and 53 are always separate, and 54 meets each of them on only 4
transition slices --- where they \emph{touch} end to end. Their regional profiles
order them: 52 at the neck, 54 at the thorax (lung 89\,\%, heart 36\,\%), 53 at
the abdomen (colon 96\,\%, small bowel 85\,\%). Alphabetical order closes the
argument, since \texttt{cervical} $<$ \texttt{lumbar} $<$ \texttt{thoracic}
reproduces exactly $52 < 53 < 54$. \texttt{54} must therefore be read as
\texttt{vertebrae\_thoracic} rather than a generic \texttt{vertebrae}.

\subsection{Conclusion: the 54 classes \emph{are} TotalSegmentator v1}

Take \texttt{class\_map["total\_v1"]} and merge the 5 lung lobes into
left/right, the 24 vertebrae into three regions, the 24 ribs into left/right and
the 5 heart parts into one:
$$104 - 5 - 24 - 24 - 5 + 8 = 54.$$
Sorted alphabetically, those 54 names land on the 54 challenge identifiers
\textbf{one for one, without exception}. The names are no longer inferred, they
are derived.

\vspace{2.0mm}

It is version 1 and not version 2, and two classes prove it. Neither appears in
the \texttt{total} task of v2: \texttt{face} was moved out to a dedicated
\texttt{face} subtask, and \texttt{pulmonary\_artery} was moved to
\texttt{heartchambers\_highres}, while a genuinely different vessel,
\texttt{pulmonary\_vein}, was added to \texttt{total} in its place. Feeding v2
labels through this table would therefore silently mislabel identifier 41 --- with
a vein, on the same anatomy.

\vspace{2.0mm}

One detail corroborates the merge itself rather than the version. Going from v1 to
v2, the authors also collapsed their five heart parts into a single \texttt{heart}
class in the \texttt{total} task: one of the four merges postulated above,
performed independently, for their own reasons.

\newpage

\section{Diagnosis: where the score is lost}
\label{sec:diagnosis}

\subsection{Every class weighs 1/54}

A class that is never predicted costs $1/54 = 0.0185$ in full, no matter how
rare it is. With 22 silent classes at the start, that is $0.41$ of the ceiling
locked away.

\vspace{2.0mm}

An oracle decomposition prices the rest: segmenting the 31 organs above 400\,px
\emph{perfectly} would still cap out at $0.6814$. In other words, no score beyond
$0.68$ is reachable without making the small organs work.

\subsection{Silent $=$ small \textbf{and} rare}

The counts below are read on the DINOv3 model, not on the baseline: by that point
the architecture change had already taken the silent classes from 22 down to 8.
Neither size nor rarity is sufficient on its own to explain which 8 survive:

\begin{center}
\begin{tabularx}{\textwidth}{lX}
\hline
\textbf{regime} & \textbf{measured} \\
\hline
rare but large & \texttt{brain} 0.987, seen 15 times; \texttt{femur\_right} 0.844, seen \textbf{9 times} \\
small but frequent & right iliac artery, 55\,px, seen 91 times: 0.115 --- predicted, but poorly \\
small \textbf{and} rare & the 8 silent classes, all under 400\,px \emph{and} under 40 slices \\
\hline
\end{tabularx}
\end{center}

\begin{sloppypar}
The double-handicap zone holds 9 classes; 8 of them are silent, and the only
survivor is \texttt{gluteus\_minimus\_left} --- a compact muscle at a
stereotyped position.
\end{sloppypar}

\subsection{Dice is the wrong instrument here}

On these classes Dice is a step function: exactly 0 until the argmax flips, then
0.3 to 0.5 at once. It cannot tell whether an intervention brought the model
\emph{closer} to the threshold, and seven of the nine organs under 200\,px scored
0.000 in every checkpoint of the project.

\vspace{2.0mm}

\noindent The fix was to read the \textbf{mean softmax probability on the organ's true
pixels} instead of the argmax. Being continuous, it separates $0.002$ from
$0.08$ where Dice writes 0 twice. Every conclusion below rests on that
instrument; without it, three of the experiments would have been declared ``no
effect'' with no way of knowing whether that was right.

\newpage

\section{What did not work}
\label{sec:failed}

\addcontentsline{toc}{subsection}{Weighting the loss by inverse organ area}
\begin{expfail}{Weighting the loss by inverse organ area}{area}
\hyp If the small classes are silent because the loss is dominated by the large
ones, re-weighting each class by the inverse of its area should bring them back.

\proto ResNet-34 U-Net, baseline recipe otherwise untouched, per-class weights
applied to the overlap term.

\meas $0.4248$, against $0.4474$ for the unweighted baseline.

\verd \rejected - and not merely neutral, \emph{below} the baseline. The weight
buys attempts on the small classes at the expense of the large ones, and the
trade loses. First sign that the small-organ failure is not a question of how the
loss is balanced.
\end{expfail}

\addcontentsline{toc}{subsection}{Raising the internal resolution}
\begin{expfail}{Raising the internal resolution}{res}
\hyp The smallest organs cover a handful of pixels. Give the network four times
as many and the argmax should be able to flip.

\proto Five paired arms (40 epochs, batch 4, ResNet-34): upsample the slice
before the network, bring the logits back down afterwards. No information is
added --- only stride.

\meas
\begin{center}
\begin{tabular}{lcc}
\hline
\textbf{arm} & \textbf{Dice} & \textbf{silent} \\
\hline
256 (control) & 0.3828 & 25 \\
384 & \textbf{0.4109} & 23 \\
512 & 0.3974 & 25 \\
256 + rare-organ resampling & 0.3903 & 24 \\
512 + resampling & 0.4023 & 23 \\
\hline
\end{tabular}
\end{center}

\verd \rejected, with a caveat that matters later. The global gain is real but
modest and \textbf{not monotone}: 384 beats 512 for half the compute. More
decisively, on the six smallest organs the probability given to the correct class
stays at \textbf{0.000 in every arm}, 512 included --- multiplying the pixels by
four changes nothing. The caveat: all five arms ran on a model carrying 23 to 25
silent classes, which is too weak to separate ``resolution does not help'' from
``nothing helps yet''. Re-run later on a model where everything else works: still
$-0.08$, and the adrenal probability falls from $0.006$ to $0.001$.
\end{expfail}

\addcontentsline{toc}{subsection}{A zoom cascade: detect, then crop and re-segment}
\begin{expfail}{A zoom cascade: detect, then crop and re-segment}{zoom}
\hyp A small organ is easier to find once the image is cropped around it.

\proto Tested directly as an \emph{upper bound}: the crop is centred on the ground
truth, so no real localisation stage could ever do better. A failure here kills
the whole family at once.

\meas
\begin{center}
\begin{tabular}{lcc}
\hline
 & \textbf{full image} & \textbf{crop 128$\to$256} \\
\hline
4 iliac vessels & prob. 0.000, Dice 0.000 & \textbf{prob. 0.28, Dice 0.22--0.26} \\
2 adrenal glands & prob. 0.000 & prob. 0.000, and falling \\
\hline
\end{tabular}
\end{center}

\verd \rejected as a pipeline. The adrenal glands do not move at all, and the four
vessels that do gain would still need a localisation stage that does not exist.

\vspace{2.0mm}
\noindent\textbf{But this is the most useful negative result of the project.}
Global upsampling $256 \to 512$ and cropping $128 \to 256$ multiply the organ's
pixel count by exactly the same factor of four. The first gives nothing; the
second takes four classes from 0 to 0.22. So what acts is not the \emph{pixel
budget} but the \textbf{narrowed field of view}: the crop removes 75\,\% of the
image, and every remaining training example is relevant to the target. A second
arm confirms it from the other side --- shrinking the decision from 55 classes to
the 12 locally plausible ones changes nothing ($0.274$ against $0.286$), so it is
not the size of the label space either.
\end{expfail}

\addcontentsline{toc}{subsection}{Per-class decision thresholds}
\begin{expfail}{Per-class decision thresholds}{thr}
\hyp The argmax assigns background as soon as its probability exceeds the
organ's. Lowering the threshold on the weak classes looks free.

\proto Fit one threshold per class, then re-score. Honest evaluation by
cross-validation, against the same fit performed and scored on all 160 validation
images.

\meas
\begin{center}
\begin{tabular}{lcccc}
\hline
\textbf{checkpoint} & \textbf{argmax} & \textbf{single thr.} & \textbf{per class} & \textbf{per class, fit on all} \\
\hline
best\_dinov3 & 0.4855 & \textbf{0.4935} & 0.4890 & 0.5077 \\
best\_resnet34 & 0.4474 & 0.4492 & \textbf{0.4587} & 0.4630 \\
pseudo\_os\_resnet34 & 0.5235 & \textbf{0.5289} & 0.5277 & 0.5416 \\
\hline
\end{tabular}
\end{center}

\verd \rejected in its per-class form. Fitting 54 thresholds on 160 images and
scoring on those same images yields $+0.018$ to $+0.022$ reproducibly;
cross-validation leaves $+0.002$ to $+0.011$. What survives is a \emph{single}
global threshold around $0.25$ --- a broad, smooth plateau and one parameter, so
it generalises. Kept in the final recipe.
\end{expfail}

\newpage

\section{What did work: external data}
\label{sec:external}

Since the 54 classes \emph{are} TotalSegmentator v1, the corresponding public
dataset (1204 CT volumes, annotated for all 104 classes, CC BY 4.0) translates
exactly. And
it brings something the challenge data does not have at all: \textbf{complete}
supervision, every organ outlined on every slice.

\subsection{Three things to imitate, or the network learns the wrong cue}

\begin{itemize}
    \item \textbf{Geometry.} $256\times256$, and above all the orientation. The
    challenge frame is rotated 90° from the usual radiological convention.
    Rather than assume it, the eight symmetries of the square were compared using
    the centroids of 12 landmark organs: the identity gives a mean error of
    $0.075$ and every other transform lands above $0.26$. The rotation added ``by
    common sense'' in the first attempt was wrong.
    \item \textbf{Plane mix.} 19\,\% of the challenge slices are not axial but
    reformats (measured: the body spans the image edge to edge). Cutting a volume
    along three planes is free, so the mix is reproduced: 81/12/7.
    \item \textbf{Intensity.} Raidium windowed every slice individually before
    quantising to 8 bits: background air lands anywhere from grey level 0 to 110,
    and one grey level is worth 4 to 32 HU depending on the slice. Each generated
    slice therefore draws its own window from the distribution measured on the
    800. Result: mean grey level 63.4 against 64.4.
\end{itemize}

\subsection{No leakage, verified}

The project's near-duplicate detector (correlation on pooled images, then mean
absolute difference at full resolution) was run between the generated slices and,
on one side the 500 test images, on the other the 800 annotated ones: no slice
below the twin threshold, and none matched to a validation image. The challenge
is therefore not drawn from the public dataset.

\subsection{What it changes for the rare classes}

\begin{center}
\begin{tabular}{llrrr}
\hline
\textbf{id} & \textbf{organ} & \textbf{challenge /800} & \textbf{external} & \textbf{ratio} \\
\hline
25 & humerus\_left & 6 & 1{,}210 & 202$\times$ \\
8 & clavicula\_right & 13 & 2{,}952 & 227$\times$ \\
2 & adrenal\_gland\_right & 18 & 1{,}780 & 99$\times$ \\
15 & gallbladder & 31 & 1{,}452 & 47$\times$ \\
\hline
\end{tabular}
\end{center}

Non-empty (image, organ) cells: $4{,}781 \to 322{,}617$, a factor of $67$.

\newpage

\section{Results}
\label{sec:results}

\begin{center}
\begin{tabular}{lccc}
\hline
\textbf{model} & \textbf{Dice (val.)} & \textbf{9 smallest} & \textbf{silent} \\
\hline
U-Net ResNet-34, baseline & 0.4474 & 0.136 & 22 \\
DINOv3 + DPT & 0.4855 & 0.165 & 8 \\
+ pseudo-labels and resampling & 0.5235 & 0.150 & 6 \\
\hline
external pretraining, 3 epochs & 0.6122 & 0.303 & 3 \\
external pretraining, 8 epochs & 0.6803 & 0.358 & \textbf{0} \\
\quad + fine-tuning on the 640 (lr $3\cdot10^{-4}$) & 0.7519 & 0.456 & 0 \\
\quad + fine-tuning on the 640 (lr $1\cdot10^{-4}$) & 0.7522 & 0.459 & 0 \\
\hline
external at $1.45$\,mm/px, pretrained & 0.6791 & 0.425 & 2 \\
\quad + fine-tuning & \textbf{0.7584} & 0.445 & 2 \\
DINOv3 + DPT, external pretrained & 0.7233 & 0.442 & 0 \\
\quad + fine-tuning & 0.7518 & 0.450 & 0 \\
\hline
ensemble, 2 ResNet-34, same scale & 0.7551 & --- & 0 \\
ensemble, 4 ResNet-34, both scales & 0.7704 & --- & 0 \\
ensemble, + DINOv3 & 0.7785 & --- & 0 \\
\textbf{ensemble, 7 members} & \textbf{0.7796} & --- & 0 \\
\hline
\end{tabular}
\end{center}

\noindent Two points deserve emphasis. First, a model pretrained purely on external data,
\textbf{having never seen a single challenge training image}, reaches $0.6122$
in three epochs - better than a ResNet-34 trained for 40 epochs on the 640
annotated slices ($0.3828$, the short-recipe control of \S\ref{sec:failed}).
Second, fine-tuning is worth $+0.072$ on its own: there remains a genuine domain
gap that imitating the geometry and the windowing does not close.

\vspace{2.0mm}

\noindent The fine-tuning learning rate, on the other hand, does not matter: $0.7519$
against $0.7522$, which is noise.

\subsection{On the leaderboard}

\begin{center}
\begin{tabular}{lccc}
\hline
\textbf{submission} & \textbf{validation} & \textbf{leaderboard} & \textbf{gap} \\
\hline
before (pseudo-labels + ensemble) & 0.5215 & 0.5226 & $+0.001$ \\
external pretraining only & 0.6122 & 0.61 & $-0.002$ \\
pretraining + fine-tuning & 0.7522 & 0.7413 & $-0.011$ \\
ensemble of 4 ResNet-34 & 0.7704 & 0.7436 & $-0.027$ \\
ensemble of 4 + DINOv3 & 0.7785 & 0.7594 & $-0.019$ \\
\textbf{+ 2 more DINOv3 arms} & \textbf{0.7796} & \textbf{0.7612} & $-0.018$ \\
\hline
\end{tabular}
\end{center}

The first row is the 4-model ensemble at $0.5215$, not the single
\texttt{pseudo\_os\_resnet34} at $0.5235$ listed above - the ensemble is what
was actually submitted at the time.

\subsection{The gap grows with how much the validation set was used}

The third and fourth rows are the cautionary ones. Early stopping picks the best
epoch out of 40 on 160 images, which is 40 draws, and costs $-0.011$. Averaging
four models each selected that way on the \emph{same} 160 images does not cancel
their bias, it compounds it - they share the ``lucky on those 160'' component,
and the gap doubles to $-0.027$. That ensemble gained $+0.018$ on validation and
$+0.002$ on the leaderboard, one eighth of what was advertised.

\vspace{2.0mm}

\noindent \textbf{And then the last row breaks the pattern, which is the useful
part.} Adding DINOv3 gained $+0.008$ on validation and $+0.016$ on the
leaderboard, twice what validation predicted, and the gap narrowed rather than
widened. The distinction is not the size of the change but its nature: a
threshold, an ensemble weighting, a post-processing sweep are knobs turned
\emph{against} 160 images and they do not transfer; an architecture, a dataset, a
training recipe are real changes and they transfer nearly in full. A member that
was not selected in the same mould also dilutes the shared bias of the others,
which is why the gap shrank.

\vspace{2.0mm}

\noindent Predicting $+0.001$ to $+0.003$ for that last submission, by
applying the one-eighth discount to what was plainly a change of model; was
the clearest analytical mistake of this project. It nearly cost the submission
that took first place.

\vspace{2.0mm}

\noindent From roughly 21\textsuperscript{st} place to
\textbf{1\textsuperscript{st}}, $0.0163$ ahead of the previous leader
($0.7449$).

\newpage

\section{The adrenal gland case}
\label{sec:adrenal}

One class remains at zero, \texttt{adrenal\_gland\_right}. The diagnosis is
worth telling because it illustrates the method.

\vspace{2.0mm}
\noindent Four hypotheses were eliminated one by one. It is not a left/right swap (the
centroids agree, $0.38/0.60$ against $0.41/0.58$). It is not a neighbouring class
stealing the pixels (the argmax says background 81\,\%, liver 18\,\%, and the
correct class ranks 6\textsuperscript{th} out of 55 at $p = 0.00005$). It is not
annotation noise (the masks are anatomically correct on both sides). It is not a
shortage of examples (1{,}780 against 18).

\vspace{2.0mm}

\noindent The decisive test: the model fails \textbf{on its own training data too}, Dice
$0.000$ over 369 external slices it has seen dozens of times. So it is neither
the domain, nor the scale, nor the quantity.

\vspace{2.0mm}

\noindent The cause is morphological. Measuring thickness (inscribed radius) rather than
area, the two adrenal glands are the \textbf{only two structures in the dataset
below 3\,px}: $2.0$ and $2.8$ px against $\geq 3.0$ for the other 52. They are
ribbons two pixels wide. But thinness alone does not explain it: the ribs are
also $3.0$\,px thick and score $0.75$ - they are bone against soft tissue. The
adrenals combine thinness \emph{and} low contrast, soft tissue against fat.

\vspace{2.0mm}

\noindent The left one unlocked by itself ($0.000 \to 0.296$) once the model became good
everywhere else. The right one, thinner by $0.8$\,px, is the last class at zero
out of 54.

\newpage

\section*{References}
\addcontentsline{toc}{section}{References}

\begin{thebibliography}{9}
\bibitem{totalseg} J. Wasserthal et al., \emph{TotalSegmentator: robust
    segmentation of 104 anatomical structures in CT images}, Radiology:
    Artificial Intelligence, 2023. Dataset:
    \url{https://zenodo.org/records/6802614} --- 1204 CT volumes, released under
    CC BY 4.0. \textbf{This work is built on that dataset and the licence
    requires the attribution.}
\bibitem{nnunet} F. Isensee et al., \emph{nnU-Net: a self-configuring method for
    deep learning-based biomedical image segmentation}, Nature Methods 18, 2021.
    The challenge's own baseline.
\bibitem{marginal} G. Shi et al., \emph{Marginal loss and exclusion loss for
    partially supervised multi-organ segmentation}, Medical Image Analysis 70,
    2021, arXiv:2007.03868.
\bibitem{dpt} R. Ranftl, A. Bochkovskiy, V. Koltun, \emph{Vision transformers
    for dense prediction}, ICCV 2021. The DPT decoder.
\bibitem{dinov3} Meta AI Research, \emph{DINOv3}, 2025. The encoder.
\end{thebibliography}

\end{document}
